\documentclass[8pt,red]{beamer}
\usepackage{amsmath,amsthm,amsfonts,amssymb,amsbsy}
\usepackage{stmaryrd} 
\usepackage{subeqnarray}

\usepackage{cases}

\usepackage{pifont}

\usepackage{fancyhdr}

\usepackage{lastpage}


\usepackage{graphicx}
\usepackage{subfloat}
\usepackage{rotating}
\usepackage{tikz}

\usetikzlibrary{arrows}
\usetikzlibrary{calc}


% \usepackage{algorithm}
% \usepackage{algorithmic}



../../../TeX/Report/macro.tex

\usepackage{wasysym}



\setbeamertemplate{theorem}[ams style]
\setbeamertemplate{theorems}[numbered]


%\input{../macro-mjean.tex}
% %hideothersubsections
% \usetheme[hideothersubsections,width=2.5cm]{Goettingen}
% %\useoutertheme[headline=empty]{miniframes}

\usetheme[]{Montpellier}
\makeatletter
\useoutertheme[height=0pt, width=0cm]{sidebar}
{\usebeamercolor{structure}}
\setbeamertemplate{sidebar canvas \beamer@sidebarside}[vertical shading][top=structure.fg!25,bottom=structure.fg!10]

%%% Local Variables: 
%%% mode: latex
%%% TeX-master: t
%%% End: 


\title{An open question : \\How to solve efficiently 3D frictional contact problem ?}
\author{Vincent Acary, Maurice Br\'emond \\ INRIA Rh\^one--Alpes, Grenoble.}
\date{Fifth symposium of the European network for nonsmooth dynamics \\ Universit\'e de Li\`ege,  September 08-09th, 2016,\\[2mm]
Joint work with
Florent Cadoux, Claude Lemar\'echal, J\'er\^ome Malick,\\ Florence Bertails--Descoubes, Gilles Daviet }

\setbeamertemplate 
{footline} 
{\quad\hfill\strut\insertsection\quad--\quad\insertframenumber/\inserttotalframenumber\strut\quad\quad} 

\graphicspath{{../figure/}}

\includeonly{%
%  introduction,
  fc3d,
  existence,
  numerics,
  comparison
}



%\newtheorem{defn}{Definition}
%\renewcommand{\thedefn}{\arabic{defn}}
% \newtheorem{thm}[defn]{Theorem}
% \newtheorem{corr}[defn]{Corollary}
% \newtheorem{ass}[defn]{Assumption}
% \newtheorem{lem}[defn]{Lemma}
% \newtheorem{rem}[defn]{Remark}
% \newtheorem{hypo}[defn]{Hypotheses}
% \newtheorem{exmp}[defn]{Example}
% \newtheorem{prop}[defn]{Proposition}
\newcommand{\diag}{\mbox{\rm diag}}
\newcommand{\co}{\overline{\mathit{co}}}
\newcommand{\rect}{\overline{\mathit{rect}}}
\newcommand{\newb}{g}

\renewcommand{\tr}[1]{\textcolor{red}{#1}}

\begin{document}

\frame{\titlepage
  \begin{center}
    \includegraphics[width=0.3\textwidth]{../RR-Inria-5.1/logo-inria.png}
  \end{center}
  $$ $$
}





%% \section{Introduction}
%% \frame
%% {
%%   \frametitle{Introduction}
%% }

% \begin{frame}
%   \frametitle{Bio.}
%   \begin{block}
%     {Team-Project BIPOP. INRIA. Centre de Grenoble Rh\^one--Alpes}
%     ``Jean Jacques Moreau's fan club''. Convex Analysis and Nonsmooth Mechanics.
%     \begin{itemize}
%     \item Scientific leader :  Bernard Brogliato
%     %\item $8$ permanents, $5$ PhD, $4$ Post-docs, $3$ Engineer,  
%     \item Nonsmooth dynamical systems :
%       Modeling, analysis, simulation and Control.
%     \item  Nonsmooth Optimization :    
%       Analysis \& algorithms. (Claude Lemar\'echal \& J\'er\^ome Malick)
%     \end{itemize}
%   \end{block}
%     \begin{block}
%       {Personal research themes}
%     \begin{itemize}
%     \item Nonsmooth Dynamical systems. Higher order Moreau's sweeping process. Complementarity systems and Filippov systems
%     \item Modeling and simulation of switched electrical circuits
%     \item Discretization method for sliding mode control and Optimal control.
%     \item Formulation and numerical solvers for Coulomb's friction and Signorini's problem. Second order cone programming.
%     \item Time--integration techniques for nonsmooth mechanical systems%: Mixed higher order schemes, Time--discontinuous Galerkin methods, Projected time--stepping schemes and generalized $\alpha$--schemes.   
%     \end{itemize}
%   \end{block}
% \end{frame}



%\section{Outline}
% \begin{frame}
%   \frametitle{Outline}
%   \begin{block}{Outline}
%     \begin{itemize}
%     \item Basic facts on piecewise linear modeling of gene regulatory network
%     \item Notion of Filippov solutions
%       \begin{itemize}
%       \item Filippov extension
%       \item Aizerman \& Pyatnitskii extension
%       \end{itemize}
%     \item Mixed Complementarity Systems formulation
%     \item Numerical time--integration
%     \item Illustrations 
%     \item Conclusions \& perspectives.
%     \end{itemize}
%   \end{block}
% \end{frame}

\frame{\tableofcontents}


%\section{Introduction}
\label{Sec:Introduction}

\frame{
  \frametitle{Contact and Friction in Mechanics}
}



%%% Local Variables: 
%%% mode: latex
%%% TeX-master: "s"
%%% End: 

%\frame{\tableofcontents}


\section{The 3D frictional contact problem}
\label{Sec:fc3d}

\subsection{Signorini condition and Coulomb's friction}

\frame{
  \frametitle{Signorini's condition and Coulomb's friction} 
  \begin{minipage}[c]{0.49\linewidth}
    \begin{tikzpicture}[ scale=2,
      axis/.style={ ->, >=stealth'},
      normal/.style={ thick, ->, >=stealth'},
      important line/.style={very thick}, 
      dashed line/.style={dashed, thin},
      every node/.style={color=black},
      soldot/.style={only marks,mark=*},
      holdot/.style={fill=white,only marks,mark=*}
      ]
      % body
      \node (BodyA) at (1,-1) {Body A};
      \fill[gray!20] (1,0) arc (0:-90:1);
      \fill[gray!20] (1,0) arc (90:180:1);
      \draw (1,0) arc (90:180:1);

      \node (BodyB) at (-1,1) {Body B};
      \draw (0,1) arc (0:-90:1);
      \fill[gray!20] (0,1) arc (90:180:1);
      \fill[gray!20] (0,1) arc (0:-90:1);

      % local frame
      \def\nlength{0.35};
      \coordinate (CA)  at  ({1.0-sqrt(2)/2.0},{-1.0+sqrt(2)/2.0});
      \node[] at  (CA) [right] {$\sf C_A$};
      \draw[holdot]  (CA) circle(0.05em);
      \draw[normal] (CA) -- ($(CA)+({-\nlength*sqrt(2)/2.0},{+\nlength*sqrt(2)/2.0 })$) node [right] {$\,\sf N$};
      \draw[normal] (CA) -- ($(CA)+({-\nlength*sqrt(2)/2.0},{-\nlength*sqrt(2)/2.0 })$) node [above] {$\sf T_1\quad$};
      \draw[dashed line] (BodyA) -- (BodyB);
      \draw[holdot] ($(CA)+({\nlength*sqrt(3)/2.0},{0.0})$) circle(0.2em);
      \node at ($(CA)+({\nlength*sqrt(3)/2.0},{0.0})$) [right]{$\sf T_2$};
      \draw[soldot] ($(CA)+({\nlength*sqrt(3)/2.0},{0.0})$) circle(0.02em);
      
      \coordinate (CB)  at  ({-1.0+sqrt(2)/2.0},{1.0-sqrt(2)/2.0});
      \node at  (CB) [above] {$\sf C_B$};
      \draw[holdot]  (CB) circle(0.05em);

      \draw[axis] (CA) -- (CB) node[midway, below left ] {$\sf g_\n$} ;

      % \draw[axis] (0,-0.4) -- (0,0.4) node(yline)[right] {$\sgn(x)$};
      % % lines
      % \draw[important line] (-0.4,-0.3) -- (0.   ,-.3);
      % \draw[important line] (0.0,0.3) --(.4,.3)  ;
      % \coordinate (O) at (0.0, 0.05);
      % \draw[fill] (O) circle (0.03em);
      % \draw (0.0,0.05) node[right]{$a$};
      % \draw (0.0,0.3) node[left]{$1$};
      % \draw (0.0,-0.3) node[right]{$-1$};
      % \draw[holdot] (0.0,0.3) circle (0.03em);
      % \draw[holdot] (0.0,-.3) circle (0.03em);
    \end{tikzpicture}
  \end{minipage}
\begin{minipage}[c]{0.49\linewidth}
    \begin{itemize}
    \item gap function $ g_{\n} = (C_B-C_A) \sf N.$
    \item reaction forces $$ r =  r_\n {\sf N} + r_\t, \quad \text{ with  } r_\n \in \RR \text{ and } r_\t \in \RR^2.$$
    \item Signorini condition at position level
      $$  0 \leq g_{\n} \perp r_\n \geq 0. $$
    \item relative velocity
      $$u =  u_\n {\sf N} + u_\t, \quad \text{ with } u_\n \in \RR \text{ and } u_\t \in \RR^2.$$
    \item Signorini condition at velocity level
        $$\left\{\begin{array}{ll}
            0 \leq u_{\n} \perp r_\n \geq 0  &\text{ if } g_{\n} \leq 0 \\
            r_{\n} =0 &\text{ otherwise}.
          \end{array}\right.$$
    \end{itemize}

  \end{minipage}
}
\frame{
  \frametitle{Signorini's condition and Coulomb's friction} 
  Let us define the Coulomb friction  cone $K$ which is chosen as the isotropic second order cone 
  \begin{equation}
    \label{eq:CoulombCone}
    K = \{r \in \RR^3 \mid \|r_\t\| \leq \mu r_n\},
  \end{equation}
  where $\mu$ is the coefficient of friction. 

  The Coulomb friction  states for the \tr{sticking case} that 
  \begin{equation}
    \label{eq:Coulom-stick}
    u_{\t} =0,\quad r \in K
  \end{equation}
  and for the \tr{sliding case} that
  \begin{equation}
    \label{eq:Coulom-slide}
    u_{\t}  \neq 0,\quad r \in \partial K, \exists\, \alpha > 0, r_\t = -\alpha u_\t.
  \end{equation}
  

}
\frame{
  \frametitle{Signorini's condition and Coulomb's friction} 
  \begin{block}
    {Disjunctive formulation of the frictional contact behavior}
    \begin{equation}
      \label{eq:contact-disjunctive}
      \left\{\begin{array}{llr}
          r = 0  &\text{ if } g_{\n} > 0  & \text{(no contact)}\\
          r = 0,  u_\n \geq 0   &\text{ if } g_{\n} \leq 0 & \text{(take--off)} \\
          r \in K, u =0 &\text{ if } g_{\n} \leq 0 & \text{(sticking)}  \\
          r \in \partial K,u _\n=0,  \exists\,\alpha > 0, u_\t = -\alpha r_\t &\text{ if } g_{\n} \leq 0 & \text{(sliding)}  \\
        \end{array}\right.
    \end{equation}
  \end{block}

  }
\frame{
  \frametitle{Signorini's condition and Coulomb's friction} 
  
  \begin{block}{SOCCP formulation}
  
    \begin{itemize}
    \item Modified relative velocity $\hat u \in \RR^3$ defined by
      \begin{equation}
        \label{eq:modified-velocity}
        \hat u = u +\mu \|u_\t\| \sf N.
      \end{equation}
    
    \item  Second-Order Cone Complementarity Problem (SOCCP) 
      \begin{equation}
        \label{eq:contact-SOCCP}
        K^\star \ni \hat u \perp r \in K
      \end{equation}
      if $ g_\n \leq 0 $ and $r=0$ otherwise.  The set $ K^\star $ is
      the dual convex cone to $K$ defined by
      \begin{equation}
        \label{eq:dual-cone}
        K^\star = \{u \in \RR^3 \mid  r^\top u \geq 0, \quad \text{for all } r \in K   \}.
      \end{equation}
      \item 
      References : \cite{Acary.Brogliato2008, Acary.ea2010},
      \cite{DeSaxce92}.

    \end{itemize}
  \end{block}

}
\frame{
  \frametitle{Signorini's condition and Coulomb's friction} 
   \begin{figure}[htbp]
  \centering
  \resizebox{!}{0.8\textheight}{\input{../figure/cone1-b.pdf_t}}
  \caption{Coulomb's friction and the modified velocity $\widehat U$. The sliding case.}
  \label{fig:CoulombFrictionSlidingDeSaxce}
\end{figure} 
}

\subsection{3D frictional contact problems}


\frame{
  \frametitle{3D frictional contact problem} 
  \begin{block}{Multiple contact notation}
    For each contact $\alpha \in \{1,\ldots n_c\}$, we have
    \begin{itemize}
    \item the local velocity : $u^\alpha \in \RR^3$, and
      $$  u = [[u^\alpha]^\top, \alpha = 1\ldots n_c]^\top$$
    \item the local  reaction vector $r^\alpha\in \RR^3$    
      $$  r = [[r^\alpha]^\top, \alpha = 1\ldots n_c]^\top$$
    \item the local  Coulomb cone    $$K^{\alpha}  = \{r^\alpha, \|r^\alpha_\t \| \leq \mu^\alpha |r^\alpha_\n| \} \subset \RR^3$$
       and the set $K$ is the cartesian product of Coulomb's friction cone at each contact, that 
      \begin{equation}
        \label{eq:CC}
        K = \prod_{\alpha=1\ldots n_c} K^{\alpha} 
      \end{equation}
      and $K^\star$ is dual.
    \end{itemize} \end{block}
  
}





\frame{
  \frametitle{3D frictional contact problems} 
  \begin{problem}[General discrete frictional contact problem]\label{prob:I}
  Given
  \begin{itemize}
    \item a symmetric positive definite matrix ${M} \in \nbR^{n \times n}$,
    \item a vector $ {f} \in \nbR^n$,
    \item a matrix  ${H} \in \nbR^{n \times m}$,
    \item a vector $w \in \RR^{m}$,
    \item a vector of coefficients of friction $\mu \in \RR^{n_c}$,
  \end{itemize}
find three vectors $ {v} \in \nbR^n$, $u\in\RR^m$ and $r\in \RR^m$, denoted by $\mathrm{FC/I}(M,H,f,w,\mu)$  such that
\begin{equation}\label{eq:soccp1}
  \begin{cases}
    M v = {H} {r} + {f} \\[2mm]
    u = H^\top v + w \\[2mm]
    \hat u = u + g(u) \\[2mm]
    K^\star \ni {\hat u} \perp r \in K
  \end{cases}
\end{equation}
with $g(u) = [[\mu^\alpha  \|u^\alpha_\t\| {\sf N}^\alpha]^\top, \alpha = 1\ldots n_c]^\top$. 
\qed
\end{problem}

}
\frame{
  \frametitle{3D frictional contact problems} 
  \begin{problem}[Reduced discrete frictional contact problem]\label{prob:II}
    Given
    \begin{itemize}
    \item a symmetric positive semi--definite  matrix ${W} \in \nbR^{m \times m}$,
    \item a vector $ {q} \in \nbR^m$,
    \item a vector $\mu \in \RR^{n_c}$ of coefficients of friction, 
    \end{itemize}
    find two vectors $u\in\RR^m$ and $r\in \RR^m$, denoted by $\mathrm{FC/II}(W,q,\mu)$  such that
    \begin{equation}\label{eq:soccp2}
      \begin{cases}
        u =Wr +q \\[2mm]
        \hat u =u + g(u) \\[2mm]
        K^\star \ni {\hat u} \perp r \in K
      \end{cases}
    \end{equation}
    with $g(u) = [[\mu^\alpha  \|u^\alpha_\t\| {\sf N}^\alpha]^\top, \alpha = 1\ldots n_c]^\top$.
    \qed
  \end{problem}
}
\frame{
  \frametitle{3D frictional contact problems} 
  This problem is really representative of a lot of applications.

  Origin of the linear relations
  $$  M v = {H} {r} + {f},\quad u = H^\top v + w$$
  \begin{itemize}
  \item Time--discretization  of the discrete dynamical mechanical system
    \begin{itemize}
    \item Event--capturing time--stepping schemes
    \item Event--detecting time--stepping schemes (event-driven)
    \end{itemize}
  \item   Time--discretization  and space discretization of the elastodynamic problem of solids
  \item   Space discretization of the quasi--static problem of solids.
  \end{itemize}
  with a possible linearization (Newton procedure.)

}

\subsection{From the mathematical programming point of view}

\frame{
  \frametitle{From the mathematical programming point of view} 
  
  
}


%%% Local Variables: 
%%% mode: latex
%%% TeX-master: "s"
%%% End: 

\frame{\tableofcontents}


\section{An existence result}
\label{Sec:existence}

\frame{
  \frametitle{An existence result. (F. Cadoux PhD)}
  Let us introduce a slack variable
  \begin{equation} \label{eq:chgVarDS}
    s^\alpha := \| u^\alpha_\t \|
  \end{equation}
  together with a new modified velocity
  $$ \widetilde u := u + \alpha s $$ $g(u) = \alpha s$
  The problem $\mathrm{FC/I}(M,H,f,w,\mu)$ can be reformulated as
  \begin{equation}\label{eq:soccp1-convex}
  \begin{cases}
    M v = {H} {r} + {f} \\[2mm]
    \widetilde  u = H^\top v + w  + \alpha s  \\[2mm]
    K^\star \ni {\hat u} \perp r \in K
  \end{cases}
\end{equation}
}
\frame{
  \frametitle{An existence result.}
  The problem~(\ref{eq:soccp1-convex}) appears to be the KKT condition of
  \begin{block}{primal problem}
    \begin{equation} \tag{$D_s$}
      \left\{
        \begin{array}{l}
          \min \quad J(v) := \frac{1}{2} v^\top M v + f^\top v \\
          H^\top v + w + \alpha s \in K^\star
        \end{array}
      \right.
    \end{equation}
  \end{block}
  \begin{block}{dual problem}
    \begin{equation} \tag{$P_s$}
      \left\{
        \begin{array}{l}
          \min \quad J_s(r) := \frac{1}{2} r^\top W r - q_s^\top r \\
          r \in K
        \end{array}
      \right.
    \end{equation}
    with $q_s = q + \alpha s $
  \end{block}
  \begin{block}{Interest}
    Two convex program \ding{220} existence of solutions under feasibility conditions.
  \end{block}
}
\frame{
  \frametitle{An existence result.}
  \begin{block}{Fixed point problem}
    Introducing
  $$ u(s) := \argmin_u (P_s) = \argmin_u (D_s) $$
  practically \alert{computable} by optimization software, and
  $$ F^i(s) := \| u^i_T(s) \|, $$
  the incremental problem becomes a fixed point problem
  $$ F(s) = s $$
\end{block}

}
\frame{
  \frametitle{An existence result.}

  \begin{block}{Assumption}
    $$ \exists v \in \R^m \ : \ H v + w \in K^* $$
  \end{block}
  Using the assumption,
  \begin{itemize}
  \item the application $F: \R_+^n \rightarrow \R_+^n$ is \alert{well-defined}, \alert{continuous} and \alert{bounded}
  \item apply Brouwer's theorem
  \end{itemize}
  \begin{theorem}
    A fixed point exists
  \end{theorem}
  This result is a variant of a previous result obtained by \cite{Pang.Klarsbring}.
}
\frame{
  \frametitle{An existence result.}
  \begin{block}{Numerical validation of the assumption}
    The assumption by solving a linear program over of a SOC.
  \end{block}
  \begin{block}{Numerical interest}
    The fixed point equation $F(s) = s$ can be tackled by
    \begin{itemize}
    \item \alert{fixed-point} iterations
      $$ s \leftarrow F(s) $$
    \item \alert{Newton} iterations
      $$s \leftarrow \mathrm{Jac}[F](s) \backslash F(s) $$
    \item Variants possible (truncated resolution of inner problem\dots)
    \end{itemize}
  \end{block}
}


%%% Local Variables: 
%%% mode: latex
%%% TeX-master: "s"
%%% End: 

\frame{\tableofcontents}


\section{Numerics}
\label{Sec:Numerics}
\subsection{VI based methods}

\frame{
  \frametitle{Numerical solution procedure. VI based methods}
  \begin{itemize}
  \item Basic fixed point iterations with projection
    $$\sf z_{k+1} \leftarrow P_{X}(z_k - \rho_k\,F(z_k))$$
  \item Extragradient method
    $$\sf z_{k+1} \leftarrow P_X(z_{k}-\rho_k\,F(P_X(z_k-\rho_k F(z_k))))$$
  \item Hyperplane projection method
  \end{itemize}
  plus self-adaptive procedure for $\rho_k$.
}


\subsection{Nonsmooth Equations based methods}
\frame{
  \frametitle{Numerical solution procedure. Nonsmooth Equations based methods}
  Nonsmooth Newton on   $G(z)=0$
   $$z_{k+1}  =  z_k -  \Phi^{-1}(z_k) (G(z_k)), \Phi(z_k) \in \partial G(z_k)$$
   \begin{itemize}
   \item Alart--Curnier Formulation
     \begin{equation}
       \label{eq:CKPS-1}
       \begin{cases}
         r_\n - P_{\RR^{n_c}_+}(r_\n - \rho_\n  u_\n) = 0, \\
         r_\t - P_{D(\mu, r_{\n,+})}(r_\t - \rho_\t u_\t   )=0,
       \end{cases}
     \end{equation}
   \item Direct normal map reformulation
     $$     r - P_{K}\left(r  - \rho (u   + g(u))\right) = 0$$
   \item  Extension of  Fischer-Burmeister function to SOCCP
     $$\phi_{\fb}(x,y) = x+y - (x^2 + y^2 )^{1/2}$$
     with Jordan product and square root
   \end{itemize}
}

\subsection{Optimization based approach}
\frame{
  \frametitle{Numerical solution procedure. Optimisation based methods}
  \begin{itemize}
  \item Successive approximation with Tresca friction (Haslinger et al.)
    \begin{equation}
      \label{eq:Haslinger-2}
      \begin{cases}
        \theta = h(r_\n) \\[2mm]
        \min\, \Frac 1 2 r^\top W r + r^\top q \\
        \begin{array}{ll}
          \text{s.t. }& r \in C(\mu,\theta)
        \end{array}
      \end{cases}
    \end{equation}
  \item Fixed on the norm of the tangential velocity (Cadoux et al.)
    \begin{equation}\label{eq:ACLM-3}
      \begin{cases}
        s =  \| u_\t \| \\[2mm]
        \min\,\Frac 1 2 r^\top W r + r^\top (q + \alpha s)  \\
        \begin{array}{ll}
          \text{s.t. } & r \in K
        \end{array}
      \end{cases}
    \end{equation}
  \item  Alternating optimization problems (Panagiotopoulos et al.)\\
    problem with fixed tangential reaction to solve normal reaction\\
    problem with fixed normal reaction to solve tangential reaction
\end{itemize}

}






%%% Local Variables: 
%%% mode: latex
%%% TeX-master: "s"
%%% End: 

\frame{\tableofcontents}


\section{First Comparisons}
\label{Sec:Comparison}

\frame{
  \frametitle{First comparisons}

}



%%% Local Variables:
%%% mode: latex
%%% TeX-master: "s"
%%% End:









\section{Conclusions \& Perspectives}
\frame{
\frametitle{Conclusions \& Perspectives}
\begin{block}{Conclusions}
  \begin{enumerate}
  \item A bunch of articles in the literature
  \item No ``Swiss--knife'' solution : choose efficiency OR robustness  
  \item Newton--based solver solves efficiently the problems but robustness issues
  \item First order iterative methods solves all the problems but very slowly  
  \item The rank of the $H$ matrix (ratio number of contacts unknows/number of d.o.f) plays an important role. 
  \end{enumerate}
\end{block}
\vskip-3mm
\begin{block}{Perspectives}
  \begin{enumerate}
  \item Develop new algorithm and compare other algorithm in the literature. \\
    (issues with standard optimization software.)
  \item Study the influence of the friction coefficient, the size of problem, the conditionning of the problem , \ldots
  \item Set up a collection of benchmarks  \ding{220} FCLIB
  \end{enumerate}
\end{block}

}

\subsection{FCLIB : a collection of discrete 3D Frictional Contact (FC) problems}

\frame{
  \frametitle{FCLIB : a collection of discrete 3D Frictional Contact (FC) problems}
  Our inspiration: MCPLIB or CUTEst
  \begin{block}
    {What is FCLIB ?}
    \begin{itemize}
    \item A open source collection of Frictional Contact (FC) problems
      stored in a specific HDF5 format 
    \item A open source light implementation of Input/Output functions
      in C Language to read and write problems (Python and Matlab coming soon)
    \end{itemize}
  \end{block}


  \begin{block}
    {Goals of the project}
    Provide a standard framework for testing available and new algorithms for solving discrete frictional contact problems share common formulations of problems in order to exchange data
\end{block}
\begin{block}
    {Call for contribution}
    \alert{\url{http://fclib.gforge.inria.fr}}
  \end{block}
}

\frame
{
\centerline{\textcolor{red}{ Thank you for your attention.}}
}

 
\def\newblock{}
{\scriptsize
\bibliographystyle{plainnat}
\bibliography{../Report/biblio,./Biblio/NonSmooth,./Biblio/Contact}
}

%\include{PWL}


\end{document}

%%% Local Variables: 
%%% mode: latex
%%% TeX-master: t
%%% End: 
